% Pandoc template for the fed_fraud documents.
% Based on proposal_gen/templates/detailed_proposal_template/detailed_1_classic.tex
% (Classic report): centered title page, table of contents, ruled section
% headings and the Geodesic house footer. Fonts are Latin Modern (lmodern).
% Two layouts: the full report (title page and contents) and, with
% -V brief, a compact one-page-brief header in the same style.
\documentclass[letterpaper,11pt]{article}
\usepackage[T1]{fontenc}
\usepackage[utf8]{inputenc}
\usepackage{lmodern}
\usepackage[letterpaper,top=0.9in,bottom=1.2in,left=1in,right=1in,footskip=0.6in]{geometry}
\usepackage[table]{xcolor}
\usepackage{geodesicbrand}
\usepackage{amsmath,amssymb}
\usepackage{enumitem}
\usepackage{array}
\usepackage{calc}
\usepackage{longtable}
\usepackage{makecell}
\usepackage[most]{tcolorbox}
\usepackage{titlesec}
\usepackage{needspace}
\usepackage{etoolbox}
\usepackage{hyperref}
\hypersetup{colorlinks=true,urlcolor=coolblack,linkcolor=coolblack,filecolor=geoblue,
  pdftitle={Partner agencies},pdfauthor={\geofirm}}
\urlstyle{same}

\setlist{itemsep=2pt,leftmargin=1.4em}
\setlength{\parindent}{0pt}
\setlength{\parskip}{6pt}
\providecommand{\tightlist}{\setlength{\itemsep}{0pt}\setlength{\parskip}{0pt}}
\setcounter{secnumdepth}{2}

% Section style: house rule under each heading
\titleformat{\section}{\Large\bfseries\color{coolblack}}{\thesection}{0.6em}{}[{\color{coolblack}\titlerule[0.5pt]}]
\titleformat{\subsection}{\large\bfseries\color{geomid}}{\thesubsection}{0.6em}{}
\titleformat{\subsubsection}{\normalsize\bfseries\color{geomid}}{\thesubsubsection}{0.6em}{}
\titlespacing*{\section}{0pt}{16pt}{8pt}

% Tables are written by latex/filters.lua (navy header, banded rows)
\setlength{\LTpre}{6pt}\setlength{\LTpost}{6pt}

\begin{document}

% TITLE PAGE
\begin{titlepage}
\newcommand{\HRule}{\rule{\linewidth}{0.3mm}}
\begin{center}
\vspace*{0.2in}
\textsc{\LARGE \geofirm}\\[0.6cm]
\includegraphics[width=30mm]{icon.png}\\[8mm]
\textsc{\Large Federal Fraud Expansion\\[0.15cm] Working Document}\\[0.5cm]
\HRule\\[0.4cm]
{\huge\bfseries Partner agencies}\\[0.3cm]
{\large Landscape and prioritization}\\[0.2cm]
\HRule\\[1cm]
\begin{minipage}{0.8\textwidth}\centering\large
Authored by:\\[0.18cm]
\geosigner, \geosignertitle: \geoemail
\end{minipage}\\[0.8cm]
{\large When: \compiletime}\\[1.2cm]
\geosurface{catenoid}{4.6cm}\\[0.9cm]
$\mathcal{G}(x)=x^x\qquad\nabla_{\dot\gamma}\dot\gamma=0$\\[3mm]
{\small\itshape \geopromise}
\end{center}
\end{titlepage}

\pagestyle{geodesic}
\tableofcontents
\newpage

\textbf{As of:} 2026-10-25. Public data only. Scores are our judgment
from the sources listed at the end; they are a starting point for
discussion, not a measured result.

\textbf{Update 2026-10-27:} the same method now covers the 33 largest
agencies by FY2026 outlays in \url{top33.pdf}. Ranks in the individual
briefs are now out of 33.

\hypertarget{bottom-line}{%
\section{Bottom line}\label{bottom-line}}

Start with \textbf{HHS (Medicare and Medicaid providers)}, then
\textbf{DoD contracting} and \textbf{USDA nutrition programs}. Keep
\textbf{SBA} running as the baseline that proves the method. HHS has the
most money at risk (about \$90B in FY2025 improper payments across
Medicare and Medicaid) and the richest public provider-level data, and
its hospice and home health fraud is the stated priority of the new
whistleblower program. DoD contracting is the closest match to the SBA
work: the same SAM.gov, FPDS and set-aside data, applied to the largest
contract budget in government. USDA's SNAP retailer and child meal
programs bring the day care and food fraud patterns the program also
names.

\hypertarget{which-agencies-are-in-scope}{%
\section{Which agencies are in
scope}\label{which-agencies-are-in-scope}}

america.gov/fraud is the intake form for DOJ's Whistleblower Incentive
Program (WHIP), launched 2026-10-20 with \$25M from Treasury and awards
of \$10K to \$30K when charges are filed. DOJ's terms page says tips
must concern federal health care offenses or fraud on programs run by
these \textbf{eight agencies}:

\begingroup\small\setlength{\tabcolsep}{4pt}\renewcommand{\arraystretch}{1.25}
\rowcolors{2}{white}{geolight}
\begin{longtable}[l]{@{}>{\raggedright\arraybackslash}p{\dimexpr 0.538\linewidth-2\tabcolsep}>{\raggedright\arraybackslash}p{\dimexpr 0.397\linewidth-2\tabcolsep}l@{}}
\rowcolor{coolblack}\leavevmode\color{white}\bfseries Agency & \leavevmode\color{white}\bfseries Example program named by DOJ & \leavevmode\color{white}\bfseries Brief \\
\endhead
Health and Human Services (incl.~CMS) & Health care benefit programs & \url{hhs.pdf} \\
War (Defense) & --- & \url{dod.pdf} \\
Small Business Administration & --- & \url{sba.pdf} \\
Agriculture & SNAP & \url{usda.pdf} \\
Veterans Affairs & --- & \url{va.pdf} \\
Housing and Urban Development & --- & \url{hud.pdf} \\
Education & Student aid fraud & \url{ed.pdf} \\
Labor & Unemployment fraud & \url{dol.pdf} \\
\end{longtable}\endgroup

DOJ (which pays the awards) and Treasury (which funds them) run the
program rather than being program agencies, so they are partners for
referrals, not targets. The site itself blocks automated reading, so
this list comes from DOJ's terms page, which the site links to. It
should still be checked against the live site. The wider task force and
DOJ's National Fraud Detection Center also name Treasury (EITC), SSA,
DHS, Interior and OMB; see ``Wave 2'' below.

\hypertarget{how-we-scored}{%
\section{How we scored}\label{how-we-scored}}

Each agency gets 1 to 5 on four factors, weighted toward what we can do
with public data alone.

\begingroup\small\setlength{\tabcolsep}{4pt}\renewcommand{\arraystretch}{1.25}
\rowcolors{2}{white}{geolight}
\begin{longtable}[l]{@{}>{\raggedright\arraybackslash}p{\dimexpr 0.283\linewidth-2\tabcolsep}l>{\raggedright\arraybackslash}p{\dimexpr 0.652\linewidth-2\tabcolsep}@{}}
\rowcolor{coolblack}\leavevmode\color{white}\bfseries Factor & \leavevmode\color{white}\bfseries Weight & \leavevmode\color{white}\bfseries What 5 means \\
\endhead
Exposure & 30\% & Tens of billions a year in improper payments or fraud-prone spend \\
Public data & 30\% & Record-level data that names payees (providers, contractors, retailers) \\
SBA method fit & 20\% & The SBA patterns (shell entities, shared identifiers, stacked awards, ineligible recipients) apply directly \\
Actionability & 20\% & A WHIP priority area, active enforcement, and a clear owner to take findings to \\
\end{longtable}\endgroup

\hypertarget{ranking}{%
\section{Ranking}\label{ranking}}

\begingroup\footnotesize\setlength{\tabcolsep}{4pt}\renewcommand{\arraystretch}{1.25}
\rowcolors{2}{white}{geolight}
\begin{longtable}[l]{@{}l>{\raggedright\arraybackslash}p{\dimexpr 0.169\linewidth-2\tabcolsep}lllll>{\raggedright\arraybackslash}p{\dimexpr 0.441\linewidth-2\tabcolsep}@{}}
\rowcolor{coolblack}\leavevmode\color{white}\bfseries Rank & \leavevmode\color{white}\bfseries Agency & \leavevmode\color{white}\bfseries Exposure & \leavevmode\color{white}\bfseries \makecell[tl]{Public\\data} & \leavevmode\color{white}\bfseries \makecell[tl]{SBA\\fit} & \leavevmode\color{white}\bfseries Action & \leavevmode\color{white}\bfseries Score & \leavevmode\color{white}\bfseries First prototype \\
\endhead
1 & HHS (CMS) & 5 & 5 & 4 & 5 & \textbf{4.8} & Hospice and home health provider risk map \\
2 & DoD & 4 & 5 & 5 & 3 & \textbf{4.3} & Small business set-aside contractor integrity screen \\
--- & SBA (baseline) & 3 & 5 & 5 & 4 & \textbf{4.2} & Extend current work to 7(a), 504 and 8(a) \\
3 & USDA & 4 & 3 & 4 & 5 & \textbf{3.9} & SNAP retailer and child meal sponsor anomalies \\
4 & VA & 3 & 4 & 4 & 3 & \textbf{3.5} & SDVOSB contract and GI Bill school screen \\
5 & HUD & 3 & 3 & 3 & 3 & \textbf{3.0} & Multifamily owner network and disaster grant screen \\
6 & Education & 2 & 3 & 3 & 3 & \textbf{2.7} & Ghost student risk by college \\
7 & Labor & 3 & 1 & 2 & 4 & \textbf{2.4} & State UI anomaly dashboard from aggregates \\
\end{longtable}\endgroup

Score = 0.3 × Exposure + 0.3 × Public data + 0.2 × SBA fit + 0.2 ×
Action.

\hypertarget{recommended-sequence}{%
\section{Recommended sequence}\label{recommended-sequence}}

\begin{enumerate}
\def\labelenumi{\arabic{enumi}.}
\tightlist
\item
  \textbf{Now: HHS and DoD.} HHS gives the biggest dollar story; DoD
  gives the fastest technical win because it reuses the SBA contractor
  pipeline as is.
\item
  \textbf{Next: USDA and VA.} USDA adds retailer and sponsor fraud; VA
  shares the contractor pipeline with DoD and adds veteran-owned
  set-asides.
\item
  \textbf{Later: HUD, Education, Labor.} Their public data is aggregate
  or thin. They are better approached once an agency pilot grants
  non-public data.
\end{enumerate}

\hypertarget{shared-building-blocks}{%
\section{Shared building blocks}\label{shared-building-blocks}}

Every brief reuses the same core, which is what makes a government-wide
program realistic:

\begin{itemize}
\tightlist
\item
  \textbf{Entity spine:} SAM.gov entity records (UEI, address, owners,
  exclusions) joined to USAspending awards and FPDS contracts.
\item
  \textbf{Exclusion and enforcement lists:} SAM.gov exclusions, HHS-OIG
  LEIE, DOJ press releases and OIG reports, used both as labels to train
  on and as watch lists.
\item
  \textbf{Common patterns:} many payees at one address, new entities
  winning large awards fast, shared officers across ``competing'' firms,
  payees tied to excluded parties, and spending far above peers.
\end{itemize}

\hypertarget{wave-2-candidates-outside-the-eight}{%
\section{Wave 2 candidates (outside the
eight)}\label{wave-2-candidates-outside-the-eight}}

\begingroup\small\setlength{\tabcolsep}{4pt}\renewcommand{\arraystretch}{1.25}
\rowcolors{2}{white}{geolight}
\begin{longtable}[l]{@{}>{\raggedright\arraybackslash}p{\dimexpr 0.142\linewidth-2\tabcolsep}>{\raggedright\arraybackslash}p{\dimexpr 0.416\linewidth-2\tabcolsep}>{\raggedright\arraybackslash}p{\dimexpr 0.442\linewidth-2\tabcolsep}@{}}
\rowcolor{coolblack}\leavevmode\color{white}\bfseries Agency & \leavevmode\color{white}\bfseries Why it matters & \leavevmode\color{white}\bfseries Public data \\
\endhead
Treasury / IRS & EITC improper payments \$21.1B (32.7\%) in FY2025 & Aggregate only; preparer-level data not public \\
SSA & SSI improper payments \$7.3B (11.6\%) in FY2025 & Aggregate only \\
DHS / FEMA & Disaster assistance fraud & OpenFEMA has program-level and some recipient data \\
\end{longtable}\endgroup

All three now have full briefs: \url{treasury.pdf}, \url{ssa.pdf} and
\url{dhs.pdf}; see \url{top33.pdf}.

\hypertarget{sources}{%
\section{Sources}\label{sources}}

\begin{itemize}
\tightlist
\item
  \href{https://www.justice.gov/fraud/whistleblower}{DOJ Whistleblower
  Rewards Program terms (lists the eight agencies)}
\item
  \href{https://www.latintimes.com/how-turn-fraudster-get-paid-immediately-inside-dojs-groundbreaking-whip-program-600382}{Latin
  Times: inside DOJ's WHIP program (launch date, award range, \$25M)}
\item
  \href{https://runtimewire.com/article/america-gov-fraud-whistleblower-ai-intake}{Runtime
  Wire: America.gov AI-assisted tip form}
\item
  \href{https://www.gao.gov/products/gao-26-108694}{GAO-26-108694:
  improper payments rose to \$186B in FY2025}
\item
  \href{https://www.crfb.org/blogs/federal-improper-payments-total-186-billion-fy-2025}{CRFB:
  FY2025 improper payments by program}
\item
  \href{https://www.gao.gov/products/GAO-24-105833}{GAO-24-105833:
  federal fraud losses estimated at \$233B to \$521B a year}
\item
  \href{../plan.pdf}{Working plan for this project}
\end{itemize}

\end{document}
